\section{The gap}

\subsection{What is established}

The literature reviewed in this chapter establishes more than might be expected of so young a field.
Agents can produce valid input for existing simulation codes from documents and descriptions, provided a harness checks every candidate and feeds errors back (Section 2.2).
They can take part in calibrating existing models, as the optimiser, as the organiser of a workflow or as a source of priors, though the best current results fall short of skilled experts on demanding real models (Section 2.3).
They can explain an existing optimisation model and answer what-if questions by having the solver do the computing (Section 2.4).
They can translate scenario narratives into inputs for a portfolio of existing energy and economic models, with an analyst approving each translation (Section 2.5).
And both the older literature on automation and recent experiments indicate that reliability depends as much on how work is divided between the language model, deterministic code and people as on the model itself (Section 2.6).

It would therefore be wrong to claim that using agents to configure, run or calibrate an existing model is new.
This thesis sits inside that body of work, not outside it.

\subsection{What is missing}

Three things are missing, in decreasing order of importance.

\subsubsection{Calibration where the parameters carry economic judgment}

Every calibration study in Section 2.3 fits physical or epidemiological parameters to observations.
More than one parameter set may fit equally well, but the candidates can be checked against measurements and physical plausibility.
Market equilibrium models such as GGM are calibrated differently in two respects.

First, the targets are partly projections.
GGM is calibrated to reproduce reference production and consumption, which for the base year come from statistics but for future years are taken from published outlooks \citep[p.~26]{egging2019global}.
A close fit to a future year therefore shows consistency with an outlook, not agreement with the world.

Second, and more importantly, the parameters adjusted are economic quantities: production costs, reference prices and the degree of market power exercised by each supplier \citep[p.~26]{egging2019global}.
None of them is observed directly, and different combinations can reproduce the same reference values.
The documentation states that which parameters to adjust is ``up to some extent arbitrary'', and that the analyst has to weigh choices using knowledge of both the market and the model (Section 1.4).
It adds that there are many defensible ways to calibrate the model and many wrong ones, and that a wrongly calibrated model may give biased results that stem from invalid parameter choices rather than from market logic \citep[p.~26]{egging2019global}.
The consequence is that two calibrations with equally good fits can embody different assumptions about how the market works, and can therefore respond differently when the model is used to study a new scenario.

Economists have discussed calibration in these terms for a long time.
The case for calibrating model economies to answer quantitative questions \citep{kydland1996computational} was met by the objection that calibrated parameters are often weakly founded, and that their uncertainty and the sensitivity of results to them should be reported routinely \citep{hansen1996empirical}.
Applied general equilibrium modelling treats calibration as choosing parameters so that the model reproduces a benchmark, with the choices themselves open to debate \citep{dawkins2001calibration}.
In hydrology, equifinality raises a similar problem for physical parameters \citep{beven2006manifesto}.
What none of the agent studies addresses is what this means when an agent does the calibrating.
A calibration of this kind cannot be judged only by its fit; it also has to be judged by whether the choices behind it are economically defensible, which is the kind of output \citet{zhu2026pai} describe as lacking a machine-readable signal of correctness.
They make the point for economic theory; it has not been examined for the calibration of an existing equilibrium model.
The OptArgus finding that a formulation can match the reference objective while being wrong \citep{li2026optargus} suggests what the risk looks like: a calibration can reproduce the reference values for reasons that would not hold up if the model were used for a different question.

% TODO: Confirm with Lukas that the base-year reference values are historical statistics (docs/ggm-model.md section 7).

\subsubsection{The middle of the workflow for energy-market models}

No study found applies agents to preparing the data for, or calibrating, a partial equilibrium or complementarity model of an energy market.
The nearest work either sits at the ends of the workflow, translating scenarios into inputs and results into explanations \citep{golden2026agentic, forootani2025re}, or concerns physical systems such as power grids, buildings and reservoirs.
\citet{golden2026agentic} come closest, and list a Global Gas Model among the models their framework can call, but they use every model with the calibration it already has.
The work that keeps such a model credible from one scenario to the next is untouched.
On its own this would be a weak claim, since novelty of application is thin ground for a thesis.
It matters because of the first point: it is precisely in this part of the workflow that the judgments described above are made.

\subsubsection{A view of the whole workflow}

Most studies automate one task, such as input generation, calibration or the explanation of results, and evaluate it in isolation.
Few ask, for a single model, which of the steps around it suit an agent, which suit ordinary code, and which should stay with a person, and why.
The task-based literature shows that such judgments have to be made task by task \citep{eloundou2024gpts, dellacqua2026navigating}, \citet{chen2025scienceagentbench} argue that agents should be assessed step by step before claims about whole workflows are made, and \citet{zhu2026human} show that the division of labour itself affects reliability.
Covering the whole workflow is therefore the appropriate scope for an exploratory thesis, not a contribution in its own right.

\subsection{How the thesis addresses the gap}

The rest of the thesis takes up these points for one model, following the sub-questions in Section 1.2.
Chapter 3 describes GGM and the workflow around it, and shows where the effort and the judgment lie.
Chapter 4 sets out in general terms what LLM agents are good and bad at, and which tasks are better done by ordinary code.
Chapter 5 brings the two together, walking through the GGM workflow step by step and asking at each step what an agent could do and where a person must stay involved; its treatment of calibration answers the first gap directly.
Chapters 6 to 8 sketch a system built on that analysis, discuss what should not be automated, and set out what remains for the master's thesis.
